\section{Conclusion and Limitations}
\label{sec:conclusion}

\method{} generates model-specific \ssn{}s that reduce harmful outputs to near zero on held-out checkpoints of two architectures while preserving task accuracy, and fingerprint-conditioned generation contributes substantially to these gains over a shared \ssn{}. For a provider maintaining many checkpoints of one backbone, the one-time meta-training is amortized across all of them.

\mysection{Limitations.}
\method{} is designed as an input-side defense against safety degradation from benign fine-tuning, rather than as a general jailbreak defense. White-box adaptive attacks on the generated \ssn{}, harms that emerge across turns, and output-side failures are outside our evaluation; checkpoint-specific weights alone do not provide security, and extending the \ssn{} to score generated text is a natural next step. \method{} assumes white-box access to activations and a backbone matching the hypernetwork, whose one-time meta-training uses safety data and is repeated per backbone family. Finally, our evaluation relies on two automated judges without human validation, measures over-refusal on XSTest and OR-Bench, evaluates baselines at a single operating point, reports seed variance for hypernetwork training, and covers 7--8B models; broadening each of these is left to future work.
